\documentclass[10pt,a4paper]{amsart}
\usepackage[margin=19mm]{geometry}
\usepackage{amsmath,amssymb,mathtools}
\usepackage[scr=rsfs,cal=euler]{mathalfa}
\usepackage{xfrac}
\usepackage[unicode,hidelinks,bookmarks=false]{hyperref}
\newcommand{\ie}{i.e.\ }
\newcommand{\cf}{cf.\ }
\newcommand{\resp}{respectively\ }
\newcommand{\Subsection}{\subsection}
\providecommand{\nobreakdash}{\nobreak\hbox{-}}
\setlength{\emergencystretch}{2em}
\multlinegap0pt
\usepackage{bm}
\usepackage{bbm}
\usepackage{dsfont}
\usepackage{xspace}
\usepackage{cancel}
\usepackage{mathtools}
\usepackage{amsthm}
\makeatletter
\@ifundefined{theorem}{\newtheorem{theorem}{Theorem}}{}
\makeatother
\title[New nonabelian Hodge graphs from twisted irregular connectio]{New nonabelian Hodge graphs from twisted irregular connections}
\author{Jean Dou\c{c}ot}
\date{Source: arXiv:2509.24861v1 (CC BY 4.0)}
\begin{document}
\maketitle

\noindent\emph{Source and licence.} New nonabelian Hodge graphs from twisted irregular connections. Version arXiv:2509.24861v1 (CC BY 4.0), distributed under the Creative Commons Attribution 4.0 International licence, as recorded in the arXiv licence metadata for this exact version and shown on the abstract page https://arxiv.org/abs/2509.24861v1. Full rights notice: licenses/arxiv-2509.24861-CC-BY-4.0.txt.\par\smallskip

\noindent\textit{Editorial scope.} Counted unit: the Theorem whose source label is \texttt{thm:characterisation\_fission\_graphs} in the article (source0.tex lines 726--728, source label \texttt{thm:characterisation\_fission\_graphs}), delivered together with the article's own complete original proof of it (source0.tex lines 730--758, one \texttt{proof} environment of 29 lines). No mathematics is written, repaired or re-derived: statement and proof are the article's own text line for line. Required context retained uncounted: none. Every definition, display and label of those lines is retained unchanged. Omitted from this package and not redistributed: 1--725, 729--729, 759--1345. Nothing inside a retained excerpt is omitted or altered: every retained body is delivered exactly as the article's lines are written, so no omission receipt of any kind accompanies this package. Citations: the retained excerpts carry no citation command at all, so no bibliography entry of the article is transcribed and no reference is redistributed. Reference adaptation: none was needed and none was made; every reference inside the delivered unit resolves inside it, so no unresolved reference or citation is produced. Printed slips kept literal: no printed word, symbol, hypothesis or number of the counted statement or of its proof was altered. Layout and class: the article's own class is \texttt{amsart} with packages including inputenc, babel, fontenc, amsmath, amsthm, amsfonts, amssymb, lmodern, mathrsfs, physics, hyperref, graphicx, geometry, enumerate; the wrapper is the lane's pinned \texttt{amsart} editorial wrapper at 10pt on A4, which restates the theorem-like environments the retained text needs and adds the article's own macro definitions that the retained text needs (none), with extra wrapper packages (bm, bbm, dsfont, xspace, cancel, mathtools). The delivered numbering is the unit's own. Assets: no retained span contains a figure environment, a graphics-inclusion command, a PDF-page inclusion, a table or a TikZ picture; no image file is copied into this package, the package declares no asset dependency and no image file is redistributed with it. Licence: version arXiv:2509.24861v1 of this article is distributed under the Creative Commons Attribution 4.0 International licence, as recorded in the arXiv licence metadata for this exact version and shown on the abstract page https://arxiv.org/abs/2509.24861v1. A full rights notice is delivered at licenses/arxiv-2509.24861-CC-BY-4.0.txt. Characters: every retained excerpt is pure ASCII, so the delivered engine, XeLaTeX with its default Latin Modern text font, reports no missing character.
\begin{theorem}\label{thm:characterisation_fission_graphs}
Let $\Gamma$ be a graph. Then it is a fission graph if and only if all its triangles are acute isosceles. 
\end{theorem}

\begin{proof}
Let $\Gamma=(N, B)$ be a fission graph. By definition, there exists an untwisted irregular class $\Theta$ (unique up to admissible deformations) such that $\Gamma=\Gamma_c(\Theta)$. The set $N$  of vertices of $\Gamma$ is the set of Stokes circles of $\Theta$. 
Let $i, j, k$ be three vertices of $\Gamma$, and let $\langle q_i\rangle$, $\langle q_j \rangle$, $\langle q_k \rangle$ be the corresponding Stokes circles. By definition of fission graphs, we have $B_{ij}=\deg(q_i-q_j)-1$, $B_{jk}=\deg(q_j-q_k)-1$, $B_{ik}=\deg(q_i-q_k)-1$. Without loss of generality we may assume $B_{ij}\leq B_{ik}\leq B_{jk}$. Now either $B_{ij}=B_{ik}=B_{jk}$, and in that case the triangle $ijk$ is ``equilateral'', and in particular acute isosceles, or $B_{ij}<B_{jk}$. In the latter case, we have $\deg(q_i-q_j)<\deg(q_j-q_k)$, which implies $\deg(q_i-q_k)=\deg((q_i-q_j)+(q_j-q_k))=\deg(q_j-q_k)$, so $B_{ik}=B_{jk}$ and the triangle $ijk$ is acute isosceles. 


Conversely, if $\Gamma=(N, B)$ is a graph such that all its triangles are acute isosceles, we can construct from it an untwisted fission tree. We proceed as follows.  Let $K-1$ be the maximal edge multiplicity appearing in $\Gamma$. We construct recursively, for $h=0, \dots, K$, an untwisted fission forest (i.e. a finite collection of untwisted fission trees) $\mathbf F_h$ with the following properties $(P_h)$. 
\begin{itemize}
\item The set of leaves of $\mathbf F_h$ is $N_0:=N$. 
\item The set of vertices of $\mathbf F_h$ is a disjoint union $N_0\sqcup N_1 \sqcup \dots \sqcup N_h$. If $x\in N_l$, $l=0, \dots h$, is a vertex of $\mathbf F_h$ we say that $x$ is of height $l$. 
\item For $l=1,\dots, h$, if $x, y\in N_l$, $x\neq y$ are two distinct vertices of height $l$, if $i\in N$ is any leaf that is a descendent of $x$ and $i\in N$ is any leaf that is a descendent of $y$, we have $B_{ij}\leq l-1$. Furthermore, if $x,y\in N$ are two leaves, and $1\leq l\leq h$, then $B_{xy}=l-1$ if and only if $x$ and $y$ have a common ancestor of height $l$ in $\mathbf F_h$. 
\end{itemize}
The forests $\mathbf F_h$ are defined as follows:


\begin{itemize}
\item $\mathbf F_0$ has set of vertices $N_0$ and no edges. 
\item Let us define $\mathbf F_1$. If $i,j\in N_0$, we write $i\sim_0 j$ if $B_{ij}=0$. This yields a well-defined equivalence relation on $N_0$, indeed it is reflexive and symmetric, and for the transitivity if $i, j,k\in N_0$ are pairwise distinct and such that $B_{ij}=B_{jk}=0$, then since the triangle $ijk$ is acute isosceles we also have $B_{ik}=0$. So we define $N_1$ to be the set of equivalence classes of $N_0$ for the relation $\sim_0$. The edges are defined as follows: for each $i\in N_0$ we draw an edge between $i$ and the vertex of $N_1$ representing its equivalence class. By construction, clearly $\mathbf F_1$ satisfies the property $(P_1)$

\item For $h\leq k$, assuming $\mathbf F_0, \dots, \mathbf F_{h-1}$ already defined and satisfying the sought after properties, let us define $\mathbf F_h$. If $i, j\in N_{h-1}$, we say that $i\sim_{h-1} j$ if for some (or equivalently any) leaf $x\in N_0$ descendent of $i$, and some (or equivalently any) leaf $y\in N_0$ descendent of $j$, we have $B_{xy}=h-1$. We claim that this yields a well-defined equivalence relation $\sim_{h-1}$ on $N_{h-1}$. 

To see that $\sim_{h-1}$ is well-defined, we have to check the above claim that `for some' descendent is equivalent to `for any' descendent. Keeping the same notations, let $x,x'$ be two distinct leaves of $i$, and assume that $B_{xy}=h-1$. Then by the induction hypothesis, $\mathbf F_{h-1}$ satisfies the property $(P_{h-1})$, which implies that $B_{xx'}\leq h-2$. Now consider the triangle $xx'y$ in $\Gamma$. It is is acute isosceles, and $B_{xx'}<B_{xy}$, so $B_{x'y}=B_{xy}=h-1$, which proves our claim. 

The reflexivity and symmetry are clear. For the transitivity, let $i, j, k\in N_{h-1}$  be pairwise distinct such that $i\sim_{h-1} j$ and $j\sim_{h-1} k$. Then if $x\in N_0$ is a leaf that is a descendent of $i$, $y\in N_0$ a leaf descendent of $j$, and $z\in N_0$ a leaf descendent of $k$, we have $B_{ij}=B_{jk}=h-1$, so since the triangle $ijk$ is acute isosceles we have $B_{ik}\leq h-1$. But we also have $B_{ik}\geq h-1$ (otherwise $y$ and $z$ would have a common ancestor of height $\leq h-2$ in $\mathbf F_{h-2}$, which would imply $i=k$), so $B_{ik}=h-1$ and $i\sim_{h-1}k$. 

We define $N_h$ to be the set of equivalences classes for $\sim_{h-1}$ in $N_{h-1}$, and define the edges of $\mathbf F_h$ by keeping those of $\mathbf F_{h-1}$, and adding edges between $N_{h-1}$ and $N_h$ as follows: if $i\in N_{h-1}$ we draw an edge between $i$ and its equivalence class in $N_h$. By construction, $\mathbf F_h$ satisfies the property $(P_h)$. 
\end{itemize}

Since $K-1$ is the maximal edge multiplicity, the forest $\mathbf F_K$ is connected (any two leaves have a common ancestor at height $\leq K-1$), i.e. is actually a tree. By construction, if $\Theta$ is any irregular class such that the part of height $\geq 1$ of its fission tree is equal to $\mathbf F_K$ with the heights of the vertices shifted by 1 (such an irregular class always exists), the corresponding fission diagram $\Gamma_c(\Theta)$ is exactly $\Gamma$, which concludes the proof.
\end{proof}

\end{document}
